% =====================================================================
% Слайд 16 — Что готово и план до защиты
% =====================================================================
\begin{frame}{Что готово и план до защиты}
  \small

  \textbf{Готово (наша работа поверх baseline):}
  \begin{itemize}\footnotesize\setlength{\itemsep}{0.2em}
    \item Архитектурная карта репозитория, слои в пакете
          \texttt{musictech/} (13 подпакетов).
    \item DTO под RL: \texttt{AlphaSummary}, \texttt{RLObservation},
          \texttt{RLAction}, \texttt{RLReward}.
    \item 9 ноутбуков с пошаговым разбором каждого компонента ---
          от формата партитуры до сборки $s_t$ из живого
          \texttt{HybridScoreFollower}.
    \item Проверено: $s_t$ собирается из текущего стека за один
          вызов, размерность 18, \emph{не зависит} от длины пьесы.
  \end{itemize}

  \textbf{План до защиты (приоритет):}
  \begin{enumerate}\footnotesize\setlength{\itemsep}{0.2em}
    \item Импортер ASAP в наш формат
          (\texttt{musictech/datasets/asap.py}) --- \rlhi{блокер}
          для всего обучения.
    \item Вынос метрик в \texttt{musictech/evaluation/}
          (alignment accuracy, onset-error, recovery latency).
    \item Среда Gymnasium $+$ симулятор солиста.
    \item BC на ASAP, потом PPO $+$ KL.
  \end{enumerate}

  \vspace{0.2em}
  \centering
  % СХЕМА 12 — Roadmap timeline
\begin{tikzpicture}[
  every node/.style={font=\scriptsize, align=center},
  done/.style={draw, very thick, circle, minimum size=0.55cm,
               inner sep=0pt, fill=black},
  next/.style={draw, thick, circle, minimum size=0.55cm,
               inner sep=0pt, fill=black!20},
  todo/.style={draw, thick, circle, minimum size=0.55cm,
               inner sep=0pt, fill=white},
]
  \draw[thick] (0, 0) -- (12, 0);

  \node[done] (a) at (0.8, 0) {};
  \node[done] (b) at (3.5, 0) {};
  \node[next] (c) at (6.2, 0) {};
  \node[todo] (d) at (8.9, 0) {};
  \node[todo] (e) at (11.6, 0) {};

  \node[above=4pt of a, text width=1.7cm] {\bfseries Baseline\\\itshape готов};
  \node[above=4pt of b, text width=1.7cm] {\bfseries Архитектура\\\itshape готова};
  \node[above=4pt of c, text width=1.7cm] {ASAP-импортер\\$+$ метрики};
  \node[above=4pt of d, text width=1.7cm] {BC $+$\\симулятор};
  \node[above=4pt of e, text width=1.7cm] {PPO + интеграция\\в живой стек};

  \foreach \x in {a, b, c, d, e}
    \node[below=4pt of \x, font=\tiny\itshape] {1--2 недели};
\end{tikzpicture}

\end{frame}
